\section{Test Program and Data}
\label{sec:method-matrix}

Switching and power transfer were checked during bring-up, followed by loaded and
no-load characterization.

\begin{table}[htbp]
\centering
\caption{Prototype test program.}
\label{tab:master-test-matrix}
\begin{tabularx}{\textwidth}{@{}p{0.25\textwidth}p{0.31\textwidth}X@{}}
\toprule
Test & Controlled condition & Main observations\\
\midrule
Bring-up & Limited input and selected load & Carrier, PA10, auxiliary rails, and basic transfer\\
Startup/shutdown & Initial voltage, load, and mode & Output rise/decay and enable transitions\\
Steady output & Fixed output setting and load & Voltage, ripple, current, and power\\
Load event & Attachment or removal of load & Voltage excursion and subsequent control response\\
No load & Open output at a selected setting & Bulk/output voltage and burst pattern\\
Power quality & Recorded AC input and output condition & Real power, PF, THD, and harmonics\\
Efficiency & Matched operating-point readings & AC input and DC output power\\
Thermal & Recorded output setting and load & Surface temperatures and location\\
Switching & Selected line and load point & Switch-node, command, and current where available\\
\bottomrule
\end{tabularx}
\end{table}

Missing setup details and peak markers limit the calculations, as discussed in
\cref{sec:method-uncertainty}. Plot lines connect measured points and are not
fitted converter characteristics.
